\section{Impact of Project on the Environment and Sustainability}
\label{sec:outro:environment}

\subsection{Environmental Impact}
Rail transport is widely recognized as one of the most energy-efficient and environmentally friendly modes of mass transportation. By modernizing interlocking with a software-based system, the project enhances the reliability and throughput of train operations, reducing delays and congestion on both rail and road networks. Every avoided collision, breakdown, or scheduling inefficiency indirectly reduces fuel wastage, unnecessary emissions from stalled locomotives, and reliance on road-based alternatives. Improved efficiency also makes railways more attractive as a sustainable alternative to long-haul road transport, contributing to national goals for reducing carbon emissions and environmental degradation.

\subsection{Sustainability Impact}
The shift toward locally developed, modular interlocking software reduces dependency on imported, hardware-intensive solutions that often come with high manufacturing, shipping, and long-term maintenance footprints. A software-driven approach allows future upgrades to be performed digitally without large-scale replacement of electromechanical equipment, thus lowering material consumption and reducing the environmental burden associated with hardware obsolescence. In the broader sense, this sustainability-oriented approach supports circularity by ensuring railway infrastructure can evolve incrementally rather than through costly, resource-intensive overhauls.

\subsection{Waste Reduction}
Accidents and failures in traditional railway systems generate both material and ecological waste—ranging from damaged rolling stock to fuel spillage and land disruption. By embedding multi-layer safety mechanisms and predictive diagnostics, the proposed system helps minimize such catastrophic waste events. Furthermore, the reliance on centralized digital logging and simulation reduces the need for extensive paper-based records and manual test procedures, indirectly cutting resource use. These incremental but cumulative measures position the system as a contributor to national efforts in waste reduction, responsible resource utilization, and sustainable transport modernization.